\documentclass[12pt]{article}
\usepackage[a4paper,margin=1in]{geometry}
\usepackage{amsmath, amssymb, booktabs, setspace}
\usepackage{enumitem}
\usepackage{graphicx}
\usepackage{siunitx}

\title{Solutions -- Practice Problem Set 4\\ \vspace{10pt} \large ECON 101 -- Macroeconomics}
\author{}
\vspace{-10pt}
\date{ \vspace{-30pt} Fall 2025}

\begin{document}
\maketitle


\begin{enumerate}[label=\textbf{\arabic*.}, leftmargin=1.5em]

\item \textbf{Solution.} In standard growth accounting, long-run growth in real GDP stems from three sources:
\begin{itemize}[leftmargin=1.2em]
  \item \emph{Capital deepening} (more physical capital per worker): Expands the productive capacity per worker but is subject to diminishing marginal returns. \textbf{Policy example (Canada):} Accelerated capital cost allowance (ACCA) or investment tax credits to lower the user cost of capital and stimulate private investment.
  \item \emph{Labour expansion and quality} (hours worked and human capital): More workers/hours raise aggregate output; higher education and skills raise output per worker. \textbf{Policy example:} Targeted immigration of high-skill workers; subsidized apprenticeships and upskilling programs.
  \item \emph{Total Factor Productivity (TFP)} (technology and efficiency): Better ideas, organization, and diffusion increase output for given inputs. \textbf{Policy example:} R\&D subsidies and stronger IP enforcement to stimulate innovation and diffusion.
\end{itemize}
Interactions matter: education and R\&D raise the return to capital; institutions (property rights, competition policy) support TFP growth.

\vspace{0.75em}

\item \emph{Given:} $Y = A K^{0.3} L^{0.7}$, with $A=5$, $K=400$, $L=100$.
\begin{enumerate}[label=(\alph*)]
\item \textbf{Level of real GDP.}
\[
Y = 5 \times 400^{0.3} \times 100^{0.7} \approx 5 \times 5.757 \times 26.334 \approx \mathbf{757.86}.
\]
(Any answer $\approx 758$ earns full credit.)

\item \textbf{Percentage increase in $Y$ if $K$ rises by 10\% (holding $L$ fixed).}\\
With a Cobb--Douglas, the elasticity of $Y$ with respect to $K$ equals the capital exponent $0.3$:
\[
\%\Delta Y \approx 0.3 \times 10\% = \mathbf{3.0\%}.
\]
Using the exact change: $(1.1)^{0.3}-1 \approx \mathbf{2.90\%}$.

\item \textbf{Why $<10\%$?}\\
Diminishing marginal returns to capital: holding $L$ and $A$ fixed, each additional unit of $K$ raises output by less than the previous one. Hence, a 10\% increase in $K$ yields only about a 3\% increase in $Y$.
\end{enumerate}

\vspace{0.75em}

\item \emph{Given:} $g_Y = 4\%$, $g_K = 2\%$, $g_L = 1\%$, factor shares $\alpha_K=0.4$, $\alpha_L=0.6$. The decomposition:
\[
g_Y \approx \alpha_K g_K + \alpha_L g_L + g_A.
\]
\begin{enumerate}[label=(\alph*)]
\item \textbf{Contributions.}
\[
\text{Capital: } 0.4 \times 2\% = \mathbf{0.8\%}, \qquad
\text{Labour: } 0.6 \times 1\% = \mathbf{0.6\%},
\]
\[
\text{TFP: } g_A = 4\% - (0.8\% + 0.6\%) = \mathbf{2.6\%}.
\]
\item \textbf{Interpretation of negative TFP.}\\
Negative TFP growth implies that, for given inputs, the economy becomes \emph{less} efficient---e.g., due to misallocation, regulatory frictions, supply disruptions, or slower diffusion of technology. Output grows more slowly than input growth would predict.
\end{enumerate}

\vspace{0.75em}

\item \textbf{Solution.}
\begin{enumerate}[label=(\alph*)]
\item \emph{R\&D subsidies:} Raise private returns to innovation $\Rightarrow$ faster idea creation/diffusion $\Rightarrow$ higher TFP \& long-run per-capita GDP.
\item \emph{Higher education:} Expands human capital, complements technology and capital, and raises the speed of adoption $\Rightarrow$ higher labour productivity and TFP.
\item \emph{Property rights/anti-corruption:} Improves contract enforcement and investment security, reduces rent-seeking, allocates resources to high-productivity uses $\Rightarrow$ higher TFP and capital formation.
\end{enumerate}

\end{enumerate}



\begin{enumerate}[label=\textbf{ \arabic*.}, start=5, leftmargin=1.5em]

\item \textbf{Solution.} In a closed economy, $AE=C+I+G$; in an open economy, $AE=C+I+G+NX$. \\
In Canada and other advanced economies, \emph{consumption} is typically the largest component. \\
\emph{Planned} investment ($I^P$) is firms' intended spending on structures, equipment, and inventories; \emph{actual} investment equals $I^P$ plus \emph{unplanned} inventory changes that occur when sales differ from expectations.

\vspace{0.75em}

\item \emph{Given:} $C=100+0.8Y_D$, $Y_D=Y-T$, $I=150$, $G=200$, $T=100$.
\begin{enumerate}[label=(\alph*)]
\item \textbf{AE function.}\\
$Y_D=Y-100 \Rightarrow C=100+0.8(Y-100)=20+0.8Y$. \\
Therefore,
\[
AE = C+I+G = (20+0.8Y)+150+200 = \mathbf{370+0.8Y}.
\]

\item \textbf{Equilibrium $Y$.}\\
Set $Y=AE$:
\[
Y=370+0.8Y \;\Rightarrow\; 0.2Y=370 \;\Rightarrow\; \boxed{Y=\;1850}.
\]

\item \textbf{MPC.} \quad $\boxed{\text{MPC}=0.8.}$

\item \textbf{Multiplier.} \quad $\displaystyle \boxed{\text{Multiplier}=\frac{1}{1-\text{MPC}}=\frac{1}{0.2}=5.}$

\item \textbf{Effect of $\Delta G=+50$.}\\
$\Delta Y = \text{multiplier}\times \Delta G = 5\times 50 = 250$. \\
New equilibrium: $\boxed{Y' = 1850+250 = 2100.}$
\end{enumerate}

\vspace{0.75em}

\item \textbf{Solution.}
\begin{enumerate}[label=(\alph*)]
\item If actual output $>$ planned $AE$, sales fall short of production $\Rightarrow$ \emph{unplanned inventory accumulation}. Firms respond by cutting production (and possibly prices), moving $Y$ toward equilibrium.
\item If actual output $<$ planned $AE$, sales exceed production $\Rightarrow$ \emph{unplanned inventory depletion}. Firms respond by raising production (and possibly prices), moving $Y$ toward equilibrium.
\end{enumerate}

\vspace{0.75em}

\item \emph{Given:} $NX=50-0.1Y$, and parameters from Question 6. \\
\textbf{Solution.} With $C=20+0.8Y$, $I=150$, $G=200$, and $NX=50-0.1Y$, we have:
\[
AE = (20+0.8Y) + 150 + 200 + (50 - 0.1Y) = \mathbf{420 + 0.7Y}.
\]
Equilibrium:
\[
Y = 420 + 0.7Y \;\Rightarrow\; 0.3Y = 420 \;\Rightarrow\; \boxed{Y=\;1400}.
\]
\emph{Multiplier intuition:} The open-economy marginal propensity to import (MPM$=0.1$) is a leakage that lowers the slope of the AE line (from $0.8$ to $0.7$ here), thus reducing the multiplier from $\frac{1}{1-0.8}=5$ to $\frac{1}{1-0.7}=\frac{1}{0.3}\approx 3.\overline{3}$. Hence, income changes are dampened when imports rise with income.
\end{enumerate}


\end{document}
